% GENERATED FILE. Edit canonical lessons, then run npm run generate:books.
% canonical-source-hash: fnv1a64:a534ddf6f41d0833
% canonical-lessons: AR-C229-rabata, AR-C229-fakka, AR-C229-saqa, AR-C229-ajana, AR-C229-masaha

\chapter{To tie, To untie, To water, To knead, To wipe}
\label{ch:ar-to-tie-to-untie-to-water-to-knead-to-wipe}
\hlchaptermodality{229}

\begin{chapteropening}
\textbf{By the end of this chapter:} \emph{I can say to tie, to untie, to water, to knead and to wipe.}

Complete the last lesson of chapter 229: I can say to tie, to untie, to water, to knead and to wipe.
\end{chapteropening}

\section[\texorpdfstring{\ar{ربط} (rabaṭa)}{ربط (rabaṭa)}]{\ar{ربط} (rabaṭa) — to tie}

\label{lesson:AR-C229-rabata}



\begin{quote}
Before the new one: say the Arabic for to sing, then the Arabic for to photograph.
\end{quote}

\subsection*{You'll want to know: \ar{ربط}}
\textbf{\ar{ربط}} — \emph{rabaṭa} — \textquotedblleft{}to tie\textquotedblright{}.

\begin{cousinweb}
\textbf{\ar{ربط}} (\emph{rabaṭa}) is built on the native Arabic root \emph{r-b-ṭ}, \textquotedblleft{}to tie\textquotedblright{} (Lane). \emph{Rabaṭa} is the plain verb, \textquotedblleft{}he tied\textquotedblright{}.
\end{cousinweb}

\begin{grammarlens}[title={What you've built}]
One more word for this chapter.
\end{grammarlens}

\subsection*{Your turn}
\begin{itemize}
\raggedright
  \item \emph{Say it:} \emph{rabaṭa}
  \item \emph{Say it:} \emph{rabaṭa}, once more
  \item \emph{Say it:} say \emph{rabaṭa}
  \item \emph{Recall:} say the Arabic for to sing, then the Arabic for to photograph, then say \emph{rabaṭa} again
\end{itemize}

\begin{tcolorbox}[breakable,colback=teal!4,colframe=teal!35!black,title={Before you move on}]
What does \ar{ربط} mean? (\textquotedblleft{}To tie\textquotedblright{}.) Say it once more.
\end{tcolorbox}
\section[\texorpdfstring{\ar{فك} (fakka)}{فك (fakka)}]{\ar{فك} (fakka) — to untie, to undo}

\label{lesson:AR-C229-fakka}



\begin{quote}
Before the new one: say the Arabic for to photograph, then the Arabic for to tie.
\end{quote}

\subsection*{You'll want to know: \ar{فك}}
\textbf{\ar{فك}} — \emph{fakka} — \textquotedblleft{}to untie, to undo\textquotedblright{}.

\begin{cousinweb}
\textbf{\ar{فك}} (\emph{fakka}) is built on the native Arabic root \emph{f-k-k}, \textquotedblleft{}to loose, to separate\textquotedblright{} (Lane). \emph{Fakka} is the plain verb, \textquotedblleft{}he undid\textquotedblright{}.
\end{cousinweb}

\begin{grammarlens}[title={What you've built}]
One more word for this chapter.
\end{grammarlens}

\subsection*{Your turn}
\begin{itemize}
\raggedright
  \item \emph{Say it:} \emph{fakka}
  \item \emph{Say it:} \emph{fakka}, once more
  \item \emph{Say it:} say \emph{fakka}
  \item \emph{Recall:} say the Arabic for to photograph, then the Arabic for to tie, then say \emph{fakka} again
\end{itemize}

\begin{tcolorbox}[breakable,colback=teal!4,colframe=teal!35!black,title={Before you move on}]
What does \ar{فك} mean? (\textquotedblleft{}To untie, to undo\textquotedblright{}.) Say it once more.
\end{tcolorbox}
\section[\texorpdfstring{\ar{سقى} (saqā)}{سقى (saqā)}]{\ar{سقى} (saqā) — to water (plants)}

\label{lesson:AR-C229-saqa}



\begin{quote}
Before the new one: say the Arabic for to tie, then the Arabic for to untie.
\end{quote}

\subsection*{You'll want to know: \ar{سقى}}
\textbf{\ar{سقى}} — \emph{saqā} — \textquotedblleft{}to water (plants)\textquotedblright{}.

\begin{cousinweb}
\textbf{\ar{سقى}} (\emph{saqā}) is built on the native Arabic root \emph{s-q-y}, \textquotedblleft{}to give to drink\textquotedblright{} (Lane). \emph{Saqā} is the plain verb, \textquotedblleft{}he watered\textquotedblright{}.
\end{cousinweb}

\begin{grammarlens}[title={What you've built}]
One more word for this chapter.
\end{grammarlens}

\subsection*{Your turn}
\begin{itemize}
\raggedright
  \item \emph{Say it:} \emph{saqā}
  \item \emph{Say it:} \emph{saqā}, once more
  \item \emph{Say it:} say \emph{saqā}
  \item \emph{Recall:} say the Arabic for to tie, then the Arabic for to untie, then say \emph{saqā} again
\end{itemize}

\begin{tcolorbox}[breakable,colback=teal!4,colframe=teal!35!black,title={Before you move on}]
What does \ar{سقى} mean? (\textquotedblleft{}To water (plants)\textquotedblright{}.) Say it once more.
\end{tcolorbox}
\section[\texorpdfstring{\ar{عجن} (ʿajana)}{عجن (ʿajana)}]{\ar{عجن} (ʿajana) — to knead}

\label{lesson:AR-C229-ajana}



\begin{quote}
Before the new one: say the Arabic for to untie, then the Arabic for to water.
\end{quote}

\subsection*{You'll want to know: \ar{عجن}}
\textbf{\ar{عجن}} — \emph{ʿajana} — \textquotedblleft{}to knead\textquotedblright{}.

\begin{cousinweb}
\textbf{\ar{عجن}} (\emph{ʿajana}) is built on the native Arabic root \emph{ʿ-j-n}, \textquotedblleft{}to knead\textquotedblright{} (Lane). \emph{ʿAjana} is the plain verb, \textquotedblleft{}he kneaded\textquotedblright{}.
\end{cousinweb}

\begin{grammarlens}[title={What you've built}]
One more word for this chapter.
\end{grammarlens}

\subsection*{Your turn}
\begin{itemize}
\raggedright
  \item \emph{Say it:} \emph{ʿajana}
  \item \emph{Say it:} \emph{ʿajana}, once more
  \item \emph{Say it:} say \emph{ʿajana}
  \item \emph{Recall:} say the Arabic for to untie, then the Arabic for to water, then say \emph{ʿajana} again
\end{itemize}

\begin{tcolorbox}[breakable,colback=teal!4,colframe=teal!35!black,title={Before you move on}]
What does \ar{عجن} mean? (\textquotedblleft{}To knead\textquotedblright{}.) Say it once more.
\end{tcolorbox}
\section[\texorpdfstring{\ar{مسح} (masaḥa)}{مسح (masaḥa)}]{\ar{مسح} (masaḥa) — to wipe}

\label{lesson:AR-C229-masaha}



\begin{quote}
Before the new one: say the Arabic for to water, then the Arabic for to knead.
\end{quote}

\subsection*{You'll want to know: \ar{مسح}}
\textbf{\ar{مسح}} — \emph{masaḥa} — \textquotedblleft{}to wipe\textquotedblright{}.

\begin{cousinweb}
\textbf{\ar{مسح}} (\emph{masaḥa}) is built on the native Arabic root \emph{m-s-ḥ}, \textquotedblleft{}to wipe\textquotedblright{} (Lane). \emph{Masaḥa} is the plain verb, \textquotedblleft{}he wiped\textquotedblright{}.
\end{cousinweb}

\begin{grammarlens}[title={What you've built}]
One more word for this chapter.
\end{grammarlens}

\subsection*{Your turn}
\begin{itemize}
\raggedright
  \item \emph{Say it:} \emph{masaḥa}
  \item \emph{Say it:} \emph{masaḥa}, once more
  \item \emph{Say it:} say \emph{masaḥa}
  \item \emph{Recall:} say the Arabic for to water, then the Arabic for to knead, then say \emph{masaḥa} again
\end{itemize}

\begin{tcolorbox}[breakable,colback=teal!4,colframe=teal!35!black,title={Before you move on}]
What does \ar{مسح} mean? (\textquotedblleft{}To wipe\textquotedblright{}.) Say it once more.
\end{tcolorbox}
